\documentclass[11pt]{article}
\usepackage{amsmath,amssymb,amsthm,hyperref}
\newtheorem{theorem}{Theorem}
\newtheorem{lemma}{Lemma}
\newtheorem{proposition}{Proposition}
\newcommand{\Prb}{\mathbb P}
\newcommand{\newsp}{\mathcal H_{\mathrm{new}}}
\title{A sharp worst-case base-size threshold for random cubic lifts}
\author{Independent research session}
\date{November 15, 2026}
\begin{document}
\maketitle

\paragraph{Status and attribution.}
This note combines a parameter refinement of the published rational trace
method with an explicit application of classical tangle theory. Historical
priority has not been established. The upper bound uses the first-order
support theorem of Chen, Garza-Vargas, Tropp and van Handel [CGVTvH]. The
covering pushforward in the lower bound is Friedman's argument [F,
Theorem 4.2]. No claim of a new spectral technique or a solution by the
mixed permutation extrema is made. The separate file
\texttt{prior/generator-free-tail-proof.md} contains the full upper-bound
derivation. The present note proves the lower bound and assembles the
sharp cubic consequence.

For an $n$-lift $G^{(n)}$ of a base $G$, define
\[
 \newsp=\left\{f:\sum_{i=1}^n f(v,i)=0\text{ for every }v\in V(G)\right\}.
\]
The old space of fibre-constant functions and its orthogonal complement
$\newsp$ both reduce adjacency. Put
\[
 \Lambda_{\mathrm{new}}(G^{(n)})=
 \|A_{G^{(n)}}|_{\newsp}\|.
\]
Permutations on distinct undirected base edges are independent and uniform.
The base is chosen deterministically before these permutations are sampled.

\begin{theorem}[Uniform upper bound]
There is an absolute constant $C$ such that, for every finite connected
loopless base multigraph $G$ with $D$ vertices, maximum degree $\Delta$,
and universal cover $T_G$,
\[
 \Prb\{\Lambda_{\mathrm{new}}(G^{(n)})\geq
        \|A_{T_G}\|+\varepsilon\}
 \leq C\frac Dn\left(\frac\Delta\varepsilon\right)^9
 \tag{1}
\]
whenever $0<\varepsilon<\Delta-\|A_{T_G}\|$. A threshold above $\Delta$
has probability zero. In particular, for each fixed $d\geq3$, every
deterministic sequence of connected $d$-regular bases with $D_n=o(n)$
satisfies
\[
 \Lambda_{\mathrm{new}}(G_n^{(n)})\leq
 2\sqrt{d-1}+o_{\Prb}(1).
 \tag{2}
\]
The constant is independent of the number of independent permutations,
which can grow with $D$.
\end{theorem}

\begin{proof}[Dependency and graph identification]
The proved parameter refinement in
\texttt{prior/generator-free-tail-proof.md} gives
\[
 \Prb\{\|P(S_1,\ldots,S_r)\|\geq\|P(u_1,\ldots,u_r)\|+
 \varepsilon\}
 \leq C\frac Dn(K/\varepsilon)^9
 \tag{3}
\]
for self-adjoint linear polynomials with $D$ by $D$ coefficients, with
no dependence on finite $r$. Here $S_i$ are the restrictions of the
permutation matrices to constants-orthogonal vectors, $u_i$ are free
regular generators, and $K$ is the full unitary norm of $P$.

Orient every base edge once and use the coefficient matrix unit from its
tail to its head. Then $P(S)$ is adjacency on $\newsp$. For arbitrary
unitary edge labels, the adjacency quadratic form is bounded by the base
adjacency form applied to the norms of its vertex blocks. Hence
$K\leq\|A_G\|\leq\Delta$. The free labelled graph has no reduced closed
walk, since distinct oriented base edges carry distinct free letters.
Each component is a universal cover of $G$, so
$\|P(u)\|=\|A_{T_G}\|$. This proves (1).
For a $d$-regular base the universal cover is the $d$-regular tree, whose
adjacency norm is $2\sqrt{d-1}$. Apply (1) at each fixed positive
$\varepsilon$ to obtain (2).
\end{proof}

Define the diamond $Q=K_4-uv$. Its vertices are $a,b,u,v$, with edges
$ab,au,av,bu,bv$. The vertices $u,v$ have degree two and will be called
ports. For $m\geq3$, let $G_m$ consist of $m$ disjoint copies $Q_j$
and the extra edges $v_j u_{j+1}$, with indices modulo $m$. This is a
connected simple cubic graph on $4m$ vertices.

Let $\alpha$ be the unique positive root greater than one of
\[
 \alpha^3-\alpha-2=0,
 \qquad t=\alpha^{-1},\qquad
 L=\alpha+2\alpha^{-1},\qquad b=(L-1)/2,
 \qquad q=2t^2.
 \tag{4}
\]
Then $\sqrt2<\alpha<2$, $0<q<1$, and
\[
 L=2.83597591908132\ldots>2\sqrt2.
 \tag{5}
\]
For an integer $R\geq0$ put
\[
 S_R=2^{R+1}+2,
 \quad M_R=2+2b^2+
          \frac{2b^2t^2(1-q^R)}{1-q},
 \quad L_R=L-\frac{2b^2tq^R}{M_R}.
 \tag{6}
\]
Thus $L_R\to L$.

\begin{theorem}[A finite lower bound on new eigenvalues]
Set
\[
 p_n=\sum_{k=1}^n\frac{(-1)^{k+1}}{k!(n)_k},
 \qquad (n)_k=n(n-1)\cdots(n-k+1).
 \tag{7}
\]
For $m\geq3$, $R\geq0$, and $n\geq2S_R$, a random $n$-lift of $G_m$
satisfies
\[
 \Prb\left\{
 \lambda_{\max}(A_{G_m^{(n)}}|_{\newsp})
       \geq L_R-\frac{6S_R}{n}\right\}
 \geq 1-(1-p_n)^m.
 \tag{8}
\]
For every fixed $\varepsilon>0$, this implies
\[
 \liminf_{n\to\infty}
 \Prb\{\Lambda_{\mathrm{new}}(G_{m_n}^{(n)})\geq L-\varepsilon\}
 \geq1-e^{-c}
 \tag{9}
\]
if $m_n/n\to c\in[0,\infty)$. If $m_n/n\to\infty$, the probability
in (9) tends to one. In particular, choosing
$0<\varepsilon<L-2\sqrt2$ disproves any uniform almost-Ramanujan
assertion that admits these bases with a number of vertices proportional
to the lift degree.
\end{theorem}

\begin{proof}[Exact core probability]
Within one block choose the spanning tree $ab,au,av$. Relabel its four
fibres so that its three edge permutations become identities. Conditional
on these three permutations, the two chord permutations remain independent
and uniform. A copy of $Q$ with one vertex in each of its four fibres is
then exactly a common fixed point of these two permutations.

Let $X$ count these copies. For $0\leq k\leq n$ its factorial moments are
\[
 \mathbb E(X)_k=(n)_k\left(\frac1{(n)_k}\right)^2
              =\frac1{(n)_k}.
 \tag{10}
\]
Inclusion-exclusion gives $\Prb(X\geq1)=p_n$. For $n\geq2$,
Bonferroni's inequalities give
\[
 \frac1n-\frac1{2n(n-1)}\leq p_n\leq\frac1n,
 \tag{11}
\]
so $p_n=1/n+O(n^{-2})$. Internal permutations in distinct blocks are
independent. Therefore the probability of at least one diamond copy is
exactly $1-(1-p_n)^m$. Connector permutations play no role in this event.
\end{proof}

\begin{proof}[Spectral witness and the new-space projection]
Complete $Q$ to an infinite cubic graph $\mathcal T(Q)$ by attaching an
infinite rooted binary branch at each port. Give $a,b$ value one, give
$u,v$ value $b$, and give every branch vertex at distance $k\geq1$ from
its port value $bt^k$. Equations (4) imply
\[
 L=\alpha+2t=t^{-1}+2t,
 \qquad L=1+2b,
 \qquad Lb=2+tb.
 \tag{12}
\]
Thus these values form a strictly positive square-summable adjacency
eigenfunction of eigenvalue $L$. Square summability follows from $q<1$.
The positive eigenfunction also bounds adjacency above by $L$, by the
weighted inequality $2xy\leq x^2 f_j/f_i+y^2 f_i/f_j$ on each edge.
Consequently $\|A_{\mathcal T(Q)}\|=L$.

Truncate the displayed function to the core and the first $R$ branch
levels. There are $S_R$ vertices. Its squared norm is $M_R$. The only
lost adjacency terms are at the last branch level, and their total
contribution to $\langle f,(L-A)f\rangle$ is $2b^2tq^R$. Its Rayleigh
quotient is therefore exactly $L_R$.

Take a nonnegative Perron vector $h$ of this finite truncated graph.
Its Perron eigenvalue $\mu_R$ is at least $L_R$. Extended by zero to
$\mathcal T(Q)$, it satisfies
\[
 A_{\mathcal T(Q)}h\geq\mu_R h
 \tag{13}
\]
pointwise. This use of the Perron vector, rather than the displayed
truncated trial vector itself, is needed for the pointwise inequality.

Suppose a finite cubic graph $H$ contains the displayed embedded diamond.
The core inclusion extends to a covering map
$\pi:\mathcal T(Q)\to H$ onto its component. This is constructed
recursively by bijectively matching the unused incident edges at each
new branch vertex. Put $g=\pi_*h$, the sum of $h$ over each inverse
image. Then $g\ne0$, $g\geq0$, and
\[
 A_Hg=\pi_*A_{\mathcal T(Q)}h\geq\mu_Rg,
 \qquad |\operatorname{supp}g|\leq S_R.
 \tag{14}
\]
Branch collisions do not invalidate this inequality. This is the covering
pushforward used in [F, Theorem 4.2].

Now let $H$ be the lift and let $g_0$ be the orthogonal projection of $g$
onto its old space. Cauchy-Schwarz on each base fibre gives
\[
 \|g_0\|^2=\frac1n\sum_{v\in V(G_m)}
     \left|\sum_i g(v,i)\right|^2
 \leq\frac{S_R}{n}\|g\|^2.
 \tag{15}
\]
Set $\delta=\|g_0\|^2/\|g\|^2\leq S_R/n\leq1/2$.
Since adjacency reduces the old and new spaces and its norm is at most
three, $g_\perp=g-g_0$ has new-space Rayleigh quotient at least
\[
 \frac{L_R-3\delta}{1-\delta}
 =L_R-\frac{(3-L_R)\delta}{1-\delta}
 \geq L_R-\frac{6S_R}{n}.
 \tag{16}
\]
In particular $g_\perp\ne0$. The variational principle proves the
deterministic spectral implication used in (8). Combine it with the exact
core probability. Finally choose $R$ fixed with $L_R>L-\varepsilon/2$,
then let $n\to\infty$ and use (11) to obtain (9).
\end{proof}

\begin{proposition}[Sharp cubic scope]
The condition $D_n=o(n)$ suffices for every deterministic connected cubic
base sequence. As a worst-case condition on the number of base vertices
alone, it cannot be relaxed to $D_n=O(n)$, even for simple bases. If
$D_n=4m_n$ and $m_n/n\to c>0$, (9) gives a failure probability at least
$1-e^{-c}$ at a fixed distance above $2\sqrt2$.
\end{proposition}

\paragraph{Limits.}
The probability bound is uniform over deterministic bases, but it is not
a simultaneous event for coefficients or bases selected after observing
the permutations. The upper bound is a direct refinement of [CGVTvH],
not a consequence of the mixed monomial extrema in the local manuscript.
No optimal epsilon exponent, full critical spectral law, edge universality,
or necessity statement for each individual base family is asserted.
No historical firstness is asserted. Numerical checks only corroborate
the formulas and are not premises of the proof.

\begin{thebibliography}{9}
\bibitem{CGVTvH}
Chi-Fang Chen, Jorge Garza-Vargas, Joel A. Tropp, and Ramon van Handel,
\emph{A new approach to strong convergence}, Annals of Mathematics 203
(2026), 555--602.
\url{https://arxiv.org/abs/2405.16026}.
Primary inputs are Lemma 5.1, the proof of Theorem 6.1, and Corollary 9.7.
The upper bound here retains a generator-independent rational denominator
and absorbs its logarithmic degree into nine derivatives.
\bibitem{F}
Joel Friedman, \emph{A proof of Alon's second eigenvalue conjecture and
related problems}, Memoirs of the AMS 195 (2008), no. 910.
\url{https://arxiv.org/abs/cs/0405020}.
Theorem 4.2 supplies the covering pushforward mechanism.
\bibitem{BC}
Charles Bordenave and Benoit Collins,
\emph{Eigenvalues of random lifts and polynomials of random permutation
matrices}, Annals of Mathematics 190 (2019), 811--875.
\url{https://arxiv.org/abs/1801.00876}.
\end{thebibliography}
\end{document}
